\clearpage
\begingroup
\fancyhead[L]{\small BEMERKUNGEN ZUM SIEBENTEN BANDE.}
\renewcommand{\headrulewidth}{0pt}
\SourceLeaf{649}
\needspace{8\baselineskip}
\pdfbookmark[1]{Bemerkungen zum siebenten Bande}{vii-volume-note}
\begin{center}\large BEMERKUNGEN ZUM SIEBENTEN BANDE.\end{center}
\begin{center}\rule{.09\linewidth}{.4pt}\end{center}

\par\noindent\hspace*{\parindent}Der vorliegende siebente Band von \textsc{Gauss}’ Werken bringt zunächst den Neudruck der Theoria motus.
Dieselbe ist bekanntlich bereits durch \textsc{Schering} im Jahre 1871 im Verlage von F. A. Perthes in einer den
Bänden I--VI der Gesammtausgabe entsprechenden Ausstattung unter Hinzufügung einiger Notizen aus dem
Nachlass neu herausgegeben worden; es schien aber trotzdem geboten, einen solchen Abdruck auch in den
vorliegenden siebenten Band der Gesammtausgabe aufzunehmen, wobei nach genauer Durchsicht noch einige
weitere Incorrectheiten der ersten Ausgabe verbessert werden konnten; auch wurden die numerischen Beispiele
einer genauen Nachrechnung unterzogen und eine grosse Reihe hier neu aufgefundener Fehler verbessert
oder, wo dies nicht möglich war, wenigstens angegeben. Ferner wurde das Werk durch eine
grosse Anzahl darauf bezüglicher neuer Notizen aus dem Nachlass bereichert.

\par\noindent\hspace*{\parindent}Ausserdem bringt der Band den bisher noch unveröffentlichten gesammten theoretisch-astronomischen
Nachlass von \textsc{Gauss}, vor allem \textsc{Gauss}’ Untersuchungen über die Störungen der Pallas, deren Bekanntmachung
von den Fachmännern seit Jahrzehnten mit Spannung erwartet wurde und die auch noch heut
grosses Interesse bieten werden.

\par\noindent\hspace*{\parindent}Zwei kleinere Aufsätze, welche bereits in Band VI aufgenommen wurden, sind wieder mitabgedruckt
worden, weil es der Zusammenhang mit den Nachlassnotizen wünschenswerth erscheinen liess, nemlich die
Aufsätze über den Zodiakus der Himmelskörper, Seite 313, und über die Tafel zur parabolischen Bewegung,
Seite 371.

\par\noindent\hspace*{\parindent}Die Bearbeitung, wie die Redaction, lag in den Händen von Herrn \textsc{Brendel} in Göttingen.

\par\noindent\hspace*{\parindent}Bei der Bearbeitung der Ceres- und der Pallasstörungen wurde eine weitgehende Ergänzung der
Nachlassnotizen durch den Bearbeiter nothwendig; es sind deshalb die Notizen, für die aus dem Nachlass
meist nur einzelne Formeln und Zahlen entnommen werden konnten, in Petit gesetzt. Ebenso wurde des
Raumes wegen der Petitsatz für den Druck der grössern Zahlentabellen gewählt, auch wenn die Zahlen aus
dem Nachlass stammen.

\par\noindent\hspace*{\parindent}Sonst sind stets, wie in den übrigen Bänden, die Einschaltungen des Bearbeiters durch eckige
Klammern [\,], die Notizen aus dem Briefwechsel u. a., die nicht von \textsc{Gauss} herrühren, in geschweifte
Klammern \(\bigl\{\bigr\}\) gesetzt.

\par\noindent\hspace*{\parindent}Weitere den Abdruck der einzelnen Notizen betreffende Hinweise findet man in den Bemerkungen
des Bearbeiters hinter jeder einzelnen Abtheilung.
\begin{flushright}\textsc{F. Klein}.\end{flushright}
\par\begin{center}\rule{.24\linewidth}{.4pt}\end{center}
\noindent\hfill81\textsuperscript{*}\par
\clearpage
\endgroup
